\documentclass{situationist-preprint}

\title{Against Arbitrary Fusions}
\msid{MET-06-b}
\msfamily{MET-06}
\msdiscipline{Metaphysics}
\mstype{Companion manuscript}
\mskeywords{special composition question; universalism; arbitrary fusions; ontological innocence; composition as identity}

\begin{document}
\maketitle

\begin{abstract}
MET-06-a argues that some things compose something exactly when they share a common fate, so that tables and organisms exist and arbitrary scattered sums do not. The universalist's best reply is that arbitrary fusions are ontologically innocent: a fusion is nothing over and above its parts, so admitting it costs nothing. We argue that innocence is a dilemma for the universalist. If fusions are nothing over and above their parts, everything true of an arbitrary fusion is true of its parts taken together, and nothing is lost by denying that there is a further object. If they are something over and above, they must earn their place, and arbitrary fusions, whose parts have no common fate, do no work that could earn it.
\end{abstract}

\section{The innocence defense}

Universalism says that any things whatever compose something: the trout and the turkey compose a trout-turkey, and a simple of a table on earth together with a simple of a star in Andromeda compose a scattered object. MET-06-a rejects this, on the ground that a composite is a locus of fate (P1, Unity, of MET-06-a), and that the members of a scattered plurality go different ways under the admissible perturbations of their environment (Assumption~3.7). The universalist replies that this objection mistakes the cost of fusions. \citet[\S 3.6]{Lewis1991} held that composition is ontologically innocent: if you are committed to the parts, you are committed to nothing further in being committed to their fusion, which is \enquote{nothing over and above} them. On the strongest version of the view, composition is identity, and the fusion is the parts \citep{Baxter1988,CotnoirBaxter2014}. An innocent object needs no fate to earn its place, and Unity is an unreasonable demand.

\section{The dilemma}

Either fusions are innocent in the sense that everything true of a fusion is true of its parts taken together, or they are not.

\paragraph{First horn.} Suppose they are. Then for every predicate true of the trout-turkey there is a plural predicate true of the trout and the turkey that says the same: the trout-turkey weighs what the trout and the turkey together weigh, is where they are, and changes when either of them changes. In that case the denial that there is a trout-turkey loses no truth. The restricted theorist of MET-06-a says that the trout and the turkey are two, and that they are, together, where they are and weigh what they weigh; the universalist adds that there is one thing which is them. If the addition is innocent, it is idle, and the restricted theory, which states the same truths without it, is the one that does not posit what does no work. Innocence makes the dispute over arbitrary fusions a dispute in which the universalist's extra sentences add nothing that the restricted theory lacks.

Strong composition as identity, moreover, does not support universalism. If the fusion of some things is them, then whether \enquote{they compose something} is true depends only on whether \enquote{they are some things}, and the special composition question receives the trivial answer: always. But the trivial answer is also available to MET-06-a, whose criterion concerns when some things compose a further object that is one in the sense that the counterfactuals of the special sciences are about it. A theory on which every plurality is, in a loose sense, one object, and on which tables are distinguished from trout-turkeys only by the behavior of their parts, has relocated the question, not answered it; and what it says about the behavior of parts is the common-fate criterion.

\paragraph{Second horn.} Suppose fusions are not innocent: a fusion is a further object, counted in addition to its parts. Then it must have a role that its parts taken together do not already play. The role of a composite in MET-06-a is to be a unified variable: the subject of counterfactuals that are stable under variations in which particular simples compose it, because its parts share a fate (MET-06-a, Section~5.1). An arbitrary fusion has no such role. There is no counterfactual about the trout-turkey that is not equivalent to a conjunction of counterfactuals about the trout and about the turkey, since nothing that happens to the one is, under the relevant dynamics, something that happens to the other. A further object with no role is posited without reason, and Unity is the statement that the reason is missing.

\section{Conclusion}

Either arbitrary fusions are innocent, in which case nothing is lost by denying them, or they are not, in which case nothing is gained by admitting them. The universalist's appeal to innocence therefore does not answer the demand of Unity, and the restriction of composition to pluralities that share a common fate stands. The argument from vagueness, the universalist's other main support \citep{Sider2001}, is answered in MET-06-a (Section~7), where the boundary of composition is shown to be sharp but unknown.

\begin{thebibliography}{9}\small
\bibitem[Baxter(1988)]{Baxter1988} Baxter, D.~L.~M. (1988). Identity in the loose and popular sense. \emph{Mind}, 97(388), 575--582.
\bibitem[Cotnoir and Baxter(2014)]{CotnoirBaxter2014} Cotnoir, A.~J. and Baxter, D.~L.~M. (eds.) (2014). \emph{Composition as Identity}. Oxford University Press.
\bibitem[Lewis(1991)]{Lewis1991} Lewis, D. (1991). \emph{Parts of Classes}. Blackwell.
\bibitem[Sider(2001)]{Sider2001} Sider, T. (2001). \emph{Four-Dimensionalism: An Ontology of Persistence and Time}. Oxford University Press.
\end{thebibliography}

\end{document}
